\documentclass[11pt,a4paper]{report}

% --- Pacchetti Fondamentali ---
\usepackage[utf8]{inputenc}
\usepackage[T1]{fontenc}
\usepackage{lmodern}
\usepackage[italian]{babel}
\usepackage{etoolbox}
\usepackage[margin=1in]{geometry}
\usepackage{hyperref}
\usepackage{xcolor}
\usepackage{listings}
\usepackage{enumitem}
\usepackage{titlesec}
\usepackage{fancyhdr}
\usepackage{cleveref}
\usepackage{graphicx}
\usepackage{longtable}
\usepackage{booktabs}
\usepackage{setspace}

\crefformat{section}{\S#2#1#3}
\crefformat{subsection}{\S#2#1#3}
\crefformat{subsubsection}{\S#2#1#3}
\crefformat{figure}{#2Figura~#1#3}
\crefformat{footnote}{#2\footnotemark[#1]#3}
\crefformat{table}{#2Tabella~#1#3}

% --- Configurazione Header e Footer ---
\pagestyle{fancy}
\fancyhf{}
\fancyhead[L]{\small Programmazione e Tecnologie Web}
\fancyhead[R]{\small Giuseppe Capasso}
\fancyfoot[C]{\thepage}
\renewcommand{\headrulewidth}{0.4pt}
\usepackage{tikz}
\usetikzlibrary{arrows.meta, positioning, shapes.geometric, shapes.misc, decorations.pathreplacing, calc, shadows, fit}

% --- Configurazione Colori e Link ---
\definecolor{primary}{RGB}{38, 66, 139}
\definecolor{codebg}{RGB}{245, 245, 245}
\definecolor{keywordcolor}{RGB}{0, 0, 255}
\definecolor{stringcolor}{RGB}{163, 21, 21}
\definecolor{commentcolor}{RGB}{0, 128, 0}

\hypersetup{
    colorlinks=true,
    linkcolor=primary,
    urlcolor=primary,
    citecolor=primary,
}

% --- Configurazione Formattazione Codice ---
\lstdefinelanguage{javascript}{
    keywords={async, await, break, case, catch, class, const, continue, debugger, default, delete, do, else, export, extends, false, finally, for, function, if, import, in, instanceof, new, null, return, super, switch, this, throw, true, try, typeof, var, void, while, with, yield, let},
    morekeywords={console, log, Error, fetch, JSON, parse, stringify, require, module, exports, on, emit, send, WebSocket},
    sensitive=true,
    morecomment=[l]{//},
    morecomment=[s]{/*}{*/},
    morestring=[b]',
    morestring=[b]",
    morestring=[b]`,
    keywordstyle=\color{primary}\bfseries,
    commentstyle=\color{commentcolor}\itshape,
    stringstyle=\color{stringcolor},
    basicstyle=\ttfamily\small
}

\lstset{
    backgroundcolor=\color{codebg},
    basicstyle=\ttfamily\small,
    breaklines=true,
    frame=single,
    rulecolor=\color{lightgray},
    captionpos=b,
    xleftmargin=1em,
    xrightmargin=1em,
    keywordstyle=\color{keywordcolor},
    commentstyle=\color{commentcolor},
    stringstyle=\color{stringcolor}
}

% --- Formattazione Titoli ---
\titleformat{\chapter}[display]
  {\normalfont\bfseries\color{primary}}{\huge Capitolo \thechapter}{1ex}{\Huge}
\titlespacing*{\chapter}{0pt}{-20pt}{40pt}

\onehalfspacing

\begin{document}

% --- Frontespizio ---
\begin{titlepage}
    \centering
    \vspace*{4cm}
    {\Huge \textbf{\color{primary} Programmazione e Tecnologie Web}} \\[1.5cm]
    {\Large \textbf{Dispense del Corso - Lezione 6}} \\[0.5cm]
    {\large Docente: Giuseppe Capasso} \\[1.5cm]
    {\normalsize \textbf{Cloud Native Development \& Cybersecurity Specialist}} \\
    {\normalsize Istituto Tecnico Superiore (ITS)} \\
    \vfill
    {\normalsize Anno Accademico 2025/2026}
\end{titlepage}

\newpage
\tableofcontents
\newpage

% ==========================================
% CAPITOLO 1
% ==========================================
\chapter{L'Evoluzione verso il Real-time}

\section{Il Modello Request-Response e i suoi Limiti}
Il World Wide Web è stato progettato originariamente per la consultazione di documenti statici tramite il protocollo HTTP. Questo modello prevede che ogni interazione sia atomica e isolata.

\begin{itemize}
    \item \textbf{Iniziativa del Client}: Il server non ha modo di "conoscere" l'indirizzo di rete del client a meno che quest'ultimo non apra una connessione.
    \item \textbf{Overhead di Connessione}: Il protocollo TCP richiede un handshake a tre vie, seguito dalla negoziazione TLS (se HTTPS), prima che il primo byte di dati HTTP possa essere scambiato.
\end{itemize}

\section{Strategie di Sincronizzazione}
Per superare la natura passiva del server, sono state implementate diverse tecniche di "polling", ognuna con un diverso impatto sulle risorse di rete.

\begin{figure}[ht]
\centering
\begin{tikzpicture}[
    node distance=1cm,
    line/.style={draw, thick, -latex},
    label/.style={font=\tiny, align=center}
]
    % --- Short Polling ---
    \node (c1) {Client};
    \node (s1) [right=2.5cm of c1] {Server};
    \draw [thick] (c1.south) -- +(0,-3.5cm);
    \draw [thick] (s1.south) -- +(0,-3.5cm);
    \node [above=0.1cm of c1, font=\bfseries\small] {Short Polling};

    \draw [line] (c1.center |- 0,-0.7) -- node[label, above] {GET /data} (s1.center |- 0,-0.7);
    \draw [line, gray, dashed] (s1.center |- 0,-1.4) -- node[label, below] {200 OK (No Data)} (c1.center |- 0,-1.4);
    
    \draw [line] (c1.center |- 0,-2.1) -- node[label, above] {GET /data} (s1.center |- 0,-2.1);
    \draw [line, primary] (s1.center |- 0,-2.8) -- node[label, below] {200 OK (New Data)} (c1.center |- 0,-2.8);

    % --- Long Polling ---
    \begin{scope}[xshift=5.5cm]
        \node (c2) {Client};
        \node (s2) [right=2.5cm of c2] {Server};
        \draw [thick] (c2.south) -- +(0,-3.5cm);
        \draw [thick] (s2.south) -- +(0,-3.5cm);
        \node [above=0.1cm of c2, font=\bfseries\small] {Long Polling};

        \draw [line] (c2.center |- 0,-0.7) -- node[label, above] {GET /data} (s2.center |- 0,-0.7);
        \draw [line, primary] (s2.center |- 0,-3.2) -- node[label, below] {200 OK (Event Triggered)} (c2.center |- 0,-3.2);
        \node [label, right=0.1cm of s2, yshift=-1.8cm, primary] {Server in attesa...};
    \end{scope}
\end{tikzpicture}
\caption{Confronto tra Short Polling e Long Polling (Lifeline corrette).}
\label{fig:polling_comparison}
\end{figure}

\section{Server-Sent Events (SSE)}
L'API SSE (\texttt{EventSource}) è una tecnologia standard che permette al server di inviare stream di dati unidirezionali al browser su HTTP. \`E ideale per scenari come quotazioni azionarie o notifiche push, dove non è richiesta un'interazione immediata del client sullo stesso canale.

\newpage

% ==========================================
% CAPITOLO 2
% ==========================================
\chapter{WebSockets e Socket.io}

\section{Il Protocollo WebSocket (RFC 6455)}
A differenza di HTTP, WebSocket fornisce un canale di comunicazione bidirezionale persistente. 

\subsection{L'Handshake Corretto}
Per evitare sovrapposizioni visive e ambiguità, analizziamo il flusso di upgrade della connessione:

\begin{figure}[ht]
\centering
\begin{tikzpicture}[
    node distance=2cm,
    box/.style={rectangle, draw, rounded corners, minimum width=2.5cm, minimum height=0.8cm, align=center, font=\small, thick},
    arrow/.style={-latex, thick}
]
    \node (client) [box, fill=blue!5] {Client\\(Browser)};
    \node (server) [box, right=5cm of client, fill=gray!10] {Server\\(Node.js)};

    % 1. Handshake Request
    \draw [arrow, primary] ([yshift=8mm]client.east) -- node[above, black, font=\tiny, align=center] {GET /chat HTTP/1.1\\Upgrade: websocket\\Connection: Upgrade} ([yshift=8mm]server.west);

    % 2. Handshake Response
    \draw [arrow, gray, dashed] ([yshift=2mm]server.west) -- node[below, black, font=\tiny, align=center] {HTTP/1.1 101 Switching Protocols\\Upgrade: websocket} ([yshift=2mm]client.east);

    % 3. Bidirectional Flow
    \draw [latex-latex, ultra thick, primary] ([yshift=-1cm]client.east) -- node[above, black, font=\small] {Stream di Frame Binari/Testo} ([yshift=-1cm]server.west);
    \draw [latex-latex, ultra thick, primary] ([yshift=-1.5cm]client.east) -- node[below, black, font=\tiny] {Bassa Latenza - Nessun Header HTTP} ([yshift=-1.5cm]server.west);

\end{tikzpicture}
\caption{Diagramma di sequenza dell'Upgrade del protocollo.}
\label{fig:ws_handshake_fixed}
\end{figure}

\section{Integrazione Client-Side (Vanilla JS)}
L'API \texttt{WebSocket} espone un'interfaccia basata su eventi per gestire il ciclo di vita della connessione.

\subsection{Eventi Fondamentali}
\begin{itemize}
    \item \textbf{onopen}: Invocato quando la connessione viene stabilita con successo. \`E il momento ideale per inviare messaggi di autenticazione o sottoscrizione.
    \item \textbf{onmessage}: Scatenato ogni volta che il server invia dati. Poich\'e i dati transitano come stringhe o buffer, \`E prassi comune utilizzare \texttt{JSON.parse()} per deserializzare gli oggetti.
    \item \textbf{onerror}: Gestisce fallimenti di connessione o errori di protocollo.
    \item \textbf{onclose}: Invocato quando la connessione viene terminata. Fondamentale per implementare logiche di riconnessione manuale.
\end{itemize}

\begin{lstlisting}[language=javascript, caption=Esempio Client WebSocket Completo]
// Inizializzazione
const socket = new WebSocket('ws://localhost:8080');

// 1. Connessione aperta
socket.onopen = () => {
    const joinMsg = { type: 'system', text: 'Utente connesso' };
    socket.send(JSON.stringify(joinMsg));
};

// 2. Ricezione dati dal server
socket.onmessage = (event) => {
    const data = JSON.parse(event.data);
    updateUI(data.user, data.text);
};

// 3. Chiusura e Errori
socket.onclose = (e) => console.log("Socket chiusa", e.code);
socket.onerror = (err) => console.error("Errore fatale", err);
\end{lstlisting}

\section{Socket.io: Oltre lo Standard}
Socket.io non è una semplice implementazione di WebSocket, ma un framework che offre funzionalità aggiuntive indispensabili in produzione:

\begin{itemize}
    \item \textbf{Fallback Automatico}: Se i WebSocket sono bloccati (es. da firewall aziendali), Socket.io ripiega su HTTP Long Polling.
    \item \textbf{Auto-reconnect}: Gestisce automaticamente il ripristino della connessione in caso di micro-interruzioni di rete.
    \item \textbf{Rooms e Namespaces}: Permette di segmentare i client in gruppi (es. stanze di chat separate) con un'API molto semplice.
\end{itemize}

\begin{lstlisting}[language=javascript, caption=Socket.io Server-side]
const io = require('socket.io')(3000);

io.on('connection', (socket) => {
    // Invio a un singolo client
    socket.emit('welcome', { id: socket.id });

    // Ricezione ed evento personalizzato
    socket.on('chat-message', (msg) => {
        // Broadcast a tutti gli altri
        socket.broadcast.emit('new-message', msg);
    });
});
\end{lstlisting}

\newpage

% ==========================================
% CAPITOLO 3
% ==========================================
\chapter{WebRTC e Multimedia}

\section{Il Problema del NAT e l'Attraversamento}
La maggior parte degli utenti internet non dispone di un indirizzo IP pubblico globale. I router domestici e aziendali utilizzano il **NAT (Network Address Translation)** per mappare più dispositivi locali su un singolo IP pubblico. Questo crea due ostacoli principali per WebRTC:
\begin{enumerate}
    \item I peer non conoscono il proprio IP pubblico e la porta mappata dal router.
    \item I router/firewall bloccano solitamente il traffico UDP in entrata non sollecitato.
\end{enumerate}

\subsection{Tipologie di NAT e Complessità di Connessione}
Non tutti i NAT si comportano allo stesso modo. La probabilità di successo di una connessione P2P diretta dipende fortemente dalla politica di filtraggio del router:

\begin{itemize}
    \item \textbf{Full Cone NAT}: Una volta mappata una porta interna su una esterna, chiunque dall'esterno può inviare pacchetti a quella porta. \`E il caso più semplice per il P2P.
    \item \textbf{Restricted Cone NAT}: Solo l'host esterno a cui il client ha inviato un pacchetto può rispondere.
    \item \textbf{Port Restricted Cone NAT}: Simile al precedente, ma restringe anche la porta sorgente dell'host esterno.
    \item \textbf{Symmetric NAT}: Ogni richiesta da una stessa porta interna verso una destinazione esterna diversa riceve una mappatura esterna diversa. Questo rende impossibile la previsione della porta e richiede obbligatoriamente un server **TURN**.
\end{itemize}

\begin{figure}[ht]
\centering
\begin{tikzpicture}[
    node distance=1.2cm,
    host/.style={draw, thick, fill=blue!5, minimum width=1.2cm, minimum height=0.8cm, font=\tiny, align=center},
    nat/.style={draw, thick, fill=red!10, chamfered rectangle, minimum width=1cm, font=\tiny, align=center},
    internet/.style={draw, thick, fill=gray!5, ellipse, minimum width=2.5cm, minimum height=1.5cm, font=\tiny},
    arrow/.style={-latex, thick}
]
    % Symmetric NAT Diagram
    \node (p1) [host] {Client Local\\192.168.1.10:5000};
    \node (n1) [nat, right=of p1] {NAT\\Symmetric};
    \node (net) [internet, right=1.5cm of n1] {Internet};
    \node (s1) [host, above right=0.5cm and 1cm of net] {Server A};
    \node (s2) [host, below right=0.5cm and 1cm of net] {Server B};

    \draw [arrow] (p1) -- (n1);
    \draw [arrow, primary] (n1.east) -- node[above, sloped, font=\tiny] {Port: 10001} (s1.west);
    \draw [arrow, primary] (n1.east) -- node[below, sloped, font=\tiny] {Port: 10002} (s2.west);
    
    \node [below=0.5cm of n1, font=\itshape\tiny, color=red] {Mappature diverse per destinazioni diverse};
\end{tikzpicture}
\caption{Visualizzazione del Symmetric NAT: la principale sfida per le connessioni P2P.}
\label{fig:nat_symmetric}
\end{figure}

\subsection{Il Framework ICE e gli Stati della Connessione}
ICE non è un protocollo a sé stante, ma un framework che combina STUN e TURN per trovare il miglior percorso di comunicazione ("Candidate").

Il ciclo di vita di una connessione WebRTC attraversa diversi stati critici monitorabili tramite \texttt{iceConnectionState}:

\begin{enumerate}
    \item \textbf{new}: L'istanza è stata creata, ma la raccolta dei candidati non è iniziata.
    \item \textbf{checking}: I peer stanno scambiando candidati e verificando la connettività.
    \item \textbf{connected}: \`E stato trovato un percorso valido per tutti i componenti.
    \item \textbf{completed}: La verifica di tutti i candidati è terminata e la connessione è stabile.
    \item \textbf{failed}: Nessun candidato ha superato i test di connettività (richiede spesso il passaggio a TURN).
    \item \textbf{disconnected/closed}: Interruzione temporanea o chiusura definitiva.
\end{enumerate}

\begin{figure}[ht]
\centering
\begin{tikzpicture}[
    node distance=1.5cm,
    state/.style={draw, thick, rounded corners, fill=primary!10, minimum width=2.2cm, font=\tiny\bfseries},
    arrow/.style={-latex, thick, gray}
]
    \node (new) [state] {NEW};
    \node (checking) [state, right=of new] {CHECKING};
    \node (connected) [state, right=of checking] {CONNECTED};
    \node (completed) [state, right=of connected] {COMPLETED};
    \node (failed) [state, below=of checking, fill=red!10] {FAILED};

    \draw [arrow] (new) -- (checking);
    \draw [arrow] (checking) -- (connected);
    \draw [arrow] (connected) -- (completed);
    \draw [arrow] (checking) -- (failed);
    \draw [arrow] (connected) -- ++(0,-1) -| (failed);
\end{tikzpicture}
\caption{Diagramma degli stati di una ICE Connection.}
\end{figure}

\section{Configurazione di una RTCPeerConnection}
L'oggetto \texttt{RTCPeerConnection} è il cuore dell'API. Gestisce la negoziazione dei codec, il superamento dei NAT e la crittografia (DTLS/SRTP).

\begin{lstlisting}[language=javascript, caption=Bootstrap di una connessione P2P]
const configuration = { 
    iceServers: [{ urls: 'stun:stun.l.google.com:19302' }] 
};
const peerConnection = new RTCPeerConnection(configuration);

// Aggiunta di stream multimediali
localStream.getTracks().forEach(track => {
    peerConnection.addTrack(track, localStream);
});

// Gestione dell'evento per ricevere stream remoti
peerConnection.ontrack = (event) => {
    const remoteVideo = document.getElementById('remoteVideo');
    remoteVideo.srcObject = event.streams[0];
};
\end{lstlisting}

\section{Conclusioni e Sfide Future}
Lo sviluppo di applicazioni real-time richiede un cambio di paradigma: non ragioniamo più per singole transazioni, ma per flussi continui di eventi. La gestione dello stato diventa distribuita e la resilienza della rete deve essere gestita esplicitamente dal programmatore.

\end{document}
